% generated from papers/XVIII-equivariant-kirchberg/paper_XVIII_equivariant_kirchberg.tex -- edit the paper, then rerun assemble.py
\chapter[The compact group advantage is illusory]{Circle actions on function field Bost--Connes systems:\\ the compact group advantage is illusory,\\ and the Weil zeta function is thermodynamic}\label{ch:XVIII}
\chapterorigin{This chapter is Paper XVIII of the series (\texttt{papers/XVIII-equivariant-kirchberg/}).}
\begin{summary}
In \S7 of Chapter~\ref{ch:XVII} we suggested that the periodicity of the modular flow over a function
field might bring the equivariant Kirchberg classification within reach, compact group actions
being better understood than flows. We now check that suggestion and withdraw it.

Takai duality converts a $\T$-action on $\cA$ into a $\Z$-action on $\cA\rtimes\T$, and $\Z$
is discrete and amenable, so the Gabe--Szab\'o equivariant Kirchberg--Phillips theorem would
indeed apply --- if $\cA\rtimes\T$ were a Kirchberg algebra. For a gauge-type action it is
the stabilized core: $\cO_n\rtimes\T$ is AF, hence stably finite. The classification theory
for Kirchberg algebras therefore does not survive the dualization. Nor does Izumi's
classification of compact group actions apply, that theory requiring the Rokhlin property,
whereas gauge actions on Cuntz algebras are approximately representable and not Rokhlin. The
advantage we conjectured does not exist.

A second obstruction is arithmetic rather than operator-algebraic. The affine boundary system
needs an infinite ring of everywhere-integral elements acting additively, and for a
projective curve a function regular at every closed point is constant:
$H^0(C,\cO_C)=\F_q$, finite. The affine construction
degenerates for $S=\cP_K$, and one must fix a point at infinity, so the classification cannot
be stated for full point sets as proposed.

Third, the equivariant $K$-theory does not contain the Weil polynomial. The Cuntz relation at
$P$ gives $\bigl(1-(qt)^{\deg P}\bigr)[1]=0$ in $K_0^\T$, so $[1]$ generates the cyclic
$R(\T)$-module $\Z[t,t^{-1}]/\bigl(1-(qt)^{d_S}\bigr)$: this records $q^{d_S}$ and nothing
more, the same information already supplied by the point spectrum
$\Sp_p(\sigma)=q^{d_S\Z}$ of Proposition 1.1 of Chapter~\ref{ch:XV}. The class group contributes
$|\Pic^0(C)(\F_q)|=P(1)$, one value of
the Weil polynomial rather than the polynomial.

What is true is the conclusion, by another route. The partition function of the Toeplitz
system is
\[
\sum_{D\in J_S}\Ns(D)^{-\beta}=\sum_{n\ge0}b_n\,q^{-\beta n}=Z_C\bigl(q^{-\beta}\bigr),
\qquad b_n=\#\{D\ge0:\deg D=n\},
\]
so the \emph{Weil zeta function is the partition function}. Knowing $\beta\mapsto Z_C(q^{-\beta})$
on $(1,\infty)$ determines the rational function $Z_C$, hence $P(T)$, the Weil numbers
$\{\alpha_j\}$, the genus, every point count $\#C(\F_{q^n})=q^n+1-\sum_j\alpha_j^n$, and
$|\Jac(C)(\F_q)|=P(1)$. The complete geometric reconstruction asked for does hold; it is
thermodynamic, not $K$-theoretic --- which is the principle of Corollary 2.4 of Chapter~\ref{ch:III}
appearing once more.
\end{summary}


\section{The affine boundary requires a point at infinity}

The framework is that of [Ch.~\ref{ch:Main}], transported to function fields as in Chapter~\ref{ch:XVII} of the
series, whose notation we keep.\begin{proposition}[Degeneration for $S=\cP_K$]\label{XVIII:prop:infinity}
Let $C/\F_q$ be a smooth projective geometrically irreducible curve and $K=\F_q(C)$. A
rational function regular at every closed point of $C$ is constant \cite{Rosen}. Now the
affine system
$\cA^{\mathrm{aff}}=C(\widehat{\cO_S})\rtimes(\cO_K\rtimes J_S)$ of Chapter~\ref{ch:II} (its \S6)
requires an \emph{infinite} additive group of everywhere-integral elements acting densely by
translations on the profinite completion --- by its Remark 6.1 the $S$-integers themselves
have poles at $S$ and do not act, which is why the additive part there is $\cO_K$. The
function-field counterpart of $\cO_K$ at $S=\cP_K$ is $H^0(C,\cO_C)=\F_q$: the additive
group is finite, its closure is a finite set, and the construction degenerates.
\end{proposition}

\begin{proof}
The first statement is that $H^0(C,\cO_C)=\F_q$ for $C$ projective and geometrically
irreducible.
\end{proof}

\begin{corollary}[The set-up must be affine]\label{XVIII:cor:affine}
One must fix a closed point $\infty\in\cP_K$ and take $A=\cO_{\cP_K\setminus\{\infty\}}$, the
ring of functions regular away from $\infty$, which is a Dedekind domain with
$\mathrm{Spec}(A)=\cP_K\setminus\{\infty\}$ infinite and dense in
$\widehat{\cO_S}$ by strong approximation, the omitted place being exactly what the theorem
requires; for $C=\mathbb P^1$ this is $\F_q[T]$.
The classification statement therefore cannot be made for full point sets
$S_i=\cP_{K_i}$: a choice of $\infty$ is part of the data, and the partition function of the
affine system is $Z_C(T)\bigl(1-T^{\deg\infty}\bigr)$ rather than $Z_C(T)$.
\end{corollary}

\section{Withdrawal: the compact group advantage does not exist}

\begin{proposition}[Takai duality lands outside the Kirchberg class]\label{XVIII:prop:takai}
Let $\alpha:\T\curvearrowright\cA$ be the gauge-type action
$\alpha_\theta(\mu_{(a,P)})=e^{i\theta\deg P}\mu_{(a,P)}$ on the affine boundary. By Takai
duality \cite{Takai}, $\alpha$
corresponds to the dual action $\hat\alpha:\Z\curvearrowright\cA\rtimes_\alpha\T$, with
$(\cA\rtimes\T)\rtimes\Z\cong\cA\otimes\mathcal K$. The group $\Z$ is discrete and amenable,
so the Gabe--Szab\'o equivariant Kirchberg--Phillips theorem \cite{GabeSzabo} applies to
$\Z$-actions \emph{on Kirchberg algebras}. However $\cA\rtimes_\alpha\T$ is stably finite,
and the Kirchberg classification does not apply to the dual system.
\end{proposition}

\begin{proof}
The action is saturated --- there are isometries of every weight in $d_S\Z$ --- so
$\cA\rtimes_\alpha\T$ is Morita equivalent, hence stably isomorphic, to the fixed-point
algebra $F=\cA^\alpha$. On $F$ the state $\tau=\lambda\circ E$, with $\lambda$ the Haar
probability of $\widehat{\cO_S}$ and $E:F\to C(\widehat{\cO_S})$ the canonical expectation,
is a faithful trace: traciality reduces to the self-similarity
$\lambda(\pi_DU+b)=\Ns(D)^{-1}\lambda(U)$ of Haar measure, the subgroup
$\pi_D\widehat{\cO_S}$ having index $\Ns(D)$, and this exactly cancels the
$\Ns(D)$ cosets appearing in $E(vxv^*)$ for a weight-$\deg D$ isometry $v$. A unital
$C^*$-algebra with a faithful tracial state is stably finite, and stable finiteness passes
to the stabilization. The model case is classical: already
$\cO_n\rtimes_\gamma\T\cong\mathcal K\otimes(\text{UHF core})$ is AF.
\end{proof}

\begin{proposition}[The action is not Rokhlin]\label{XVIII:prop:notrokhlin}
Izumi's classification of finite and compact group actions on Kirchberg algebras requires the
Rokhlin property \cite{Izumi}. A Rokhlin action of $\T$ on a unital $\cA$ forces $\cA$ to
absorb the fixed-point algebra in a strong sense; the restrictions of gauge actions on Cuntz
and Cuntz--Pimsner algebras to finite cyclic subgroups are instead \emph{approximately
representable}, the notion dual to Rokhlin \cite{Izumi}, and the circle action itself does
not have the Rokhlin property. So that route is unavailable as well.
\end{proposition}

\begin{remark}[Withdrawal of a suggestion]\label{XVIII:rem:withdraw}
We suggested there that periodicity of $\sigma$ ``may bring Remark 5.1 of Chapter~\ref{ch:XIII} within
reach'', on the ground that the equivariant classification for compact groups is better
developed than for $\R$. That is true as a general statement about the literature and false
as an inference here, for the two reasons above: dualizing a circle action moves the problem
to a stably finite algebra, and the specific action fails the Rokhlin hypothesis under which
the compact group theory is available. We withdraw the suggestion. The caution of
Remark 5.1 of Chapter~\ref{ch:XIII}, that flows on Kirchberg algebras are not classified, stands and
is not circumvented by compactness.
\end{remark}

\section{What the equivariant $K$-theory records}

\begin{proposition}[The cyclic module]\label{XVIII:prop:KT}
In $K_0^\T(\partial\cA^{\mathrm{aff}})$, the Cuntz relation at $P$ --- the sum over the
$\Ns P=q^{\deg P}$ additive cosets of the range projections, each of weight $t^{\deg P}$ ---
gives
\[
\bigl(1-(qt)^{\deg P}\bigr)\,[1]=0\qquad(P\in S).
\]
Since $x=qt$ becomes a unit modulo each relation $1-x^{\deg P}$, the exponents combine by
B\'ezout: $\gcd_i\bigl(1-x^{n_i}\bigr)=1-x^{\gcd_in_i}$ as ideals of $\Z[t,t^{-1}]$, and the
class $[1]$ generates the cyclic $R(\T)=\Z[t,t^{-1}]$-module
\[
R(\T)\cdot[1]\ \cong\ \Z[t,t^{-1}]\big/\bigl(1-(qt)^{d_S}\bigr),\qquad d_S=\gcd_{P\in S}\deg P .
\]
\end{proposition}

\begin{corollary}[No genus]\label{XVIII:cor:nogenus}
The module of Proposition~\ref{XVIII:prop:KT} records the single quantity $q^{d_S}$, which is
exactly the content of the point spectrum $\Sp_p(\sigma)=q^{d_S\Z}$ already available from
$(\cA,\sigma)$ by Proposition 1.1 of Chapter~\ref{ch:XV}. It contains no information about the genus, the
Weil numbers, or the point counts.
\end{corollary}

\begin{remark}[What $K$-theory does see]\label{XVIII:rem:whatKsees}
By Theorem 3.3 of Chapter~\ref{ch:II} the class of the unit in $K_0$ of the boundary records the
narrow class number; the counterpart of that quantity here is
$|\Pic^0(C)(\F_q)|=|\Jac(C)(\F_q)|$, of order $P(1)$. So $K$-theory does see
one value of the Weil polynomial, at $T=1$. It does not see the polynomial: the $2g$
numbers $\alpha_j$ are not determined by $P(1)$, and no mechanism in the six-term sequence
produces the higher point counts.
\end{remark}

\section{The Weil zeta function is the partition function}

\begin{theorem}[$Z_C$ as a thermodynamic quantity]\label{XVIII:thm:partition}
Let $b_n=\#\{D\in\Div_+(C):\deg D=n\}$. Then, by the definition of the Weil zeta function,
\[
Z_C(T)=\sum_{n\ge0}b_nT^n=\frac{P(T)}{(1-T)(1-qT)},\qquad P(T)=\prod_{j=1}^{2g}(1-\alpha_jT),
\]
and the partition function of the Toeplitz system $(\cA_{C,\cP_K},\sigma)$ is
\begin{equation}\label{XVIII:eq:partition}
\sum_{D\in J_{\cP_K}}\Ns(D)^{-\beta}=\sum_{n\ge0}b_n\,q^{-\beta n}=Z_C\bigl(q^{-\beta}\bigr),
\qquad\beta>1 .
\end{equation}
Equivalently, the $\T$-action on the Fock space $\ell^2(J_S)$ has weight-$n$ multiplicity
$b_n$, and $Z_C$ is its graded dimension.
\end{theorem}

\begin{proof}
$\Ns(D)=q^{\deg D}$ and $J_{\cP_K}=\Div_+(C)$; group the sum by degree. Convergence for
$\beta>1$ is Corollary 2.2 of Chapter~\ref{ch:XVII}.
\end{proof}

\begin{theorem}[Complete geometric reconstruction]\label{XVIII:thm:reconstruct}
The function $\beta\mapsto Z_C(q^{-\beta})$ on $(1,\infty)$ determines the rational function
$Z_C$, hence:
\begin{enumerate}
\item the Weil polynomial $P(T)$ and the Weil numbers $\{\alpha_1,\dots,\alpha_{2g}\}$;
\item the genus $g=\tfrac12\deg P$;
\item every point count $\#C(\F_{q^n})=q^n+1-\sum_{j=1}^{2g}\alpha_j^n$;
\item the order $|\Jac(C)(\F_q)|=P(1)$ of the group of rational points of the Jacobian.
\end{enumerate}
Consequently the equivariant $C^*$-dynamical system $(\cA_{C,\cP_K},\T,\sigma)$ is a complete
invariant of the Weil zeta function of $C$.
\end{theorem}

\begin{proof}
A rational function is determined by its values on an interval; $q$ is recovered from the
point spectrum $\Sp_p(\sigma)=q^{d_S\Z}$ (Proposition 1.1 of Chapter~\ref{ch:XV}, computed for these
systems in Remark 4.3 of Chapter~\ref{ch:XVII}), with $d_S=1$ for $S=\cP_K$ by F.~K. Schmidt
(Proposition 1.1 of Chapter~\ref{ch:XVII}), and then
$P(T)=Z_C(T)(1-T)(1-qT)$. The remaining items are the Weil conjectures for curves.
\end{proof}

\begin{example}[Numerical illustration]\label{XVIII:ex:elliptic}
For genus one, $P(T)=1-aT+qT^2$:
\[
\begin{array}{lccl}
(q,a) & b_0,\dots,b_4 & \#C(\F_{q^n}),\ n\le4 & |\Jac(\F_q)|\\\hline
(4,1) & 1,4,20,84,340 & 4,\ 24,\ 76,\ 240 & 4\\
(5,-2) & 1,8,48,248,1248 & 8,\ 32,\ 104,\ 640 & 8\\
(9,4) & 1,6,60,546,4920 & 6,\ 84,\ 774,\ 6720 & 6
\end{array}
\]
The $b_n$ are the multiplicities in the partition function \eqref{XVIII:eq:partition}; the point
counts and the Jacobian order follow from them.
\end{example}

\begin{remark}[The principle, again]\label{XVIII:rem:principle}
Theorem~\ref{XVIII:thm:reconstruct} delivers the conclusion asked for while
Corollary~\ref{XVIII:cor:nogenus} shows the proposed route cannot. This is the pattern established
in Corollary 2.4 of Chapter~\ref{ch:III}: the arithmetic of these systems lives in the equilibrium data
--- the measure class, the KMS simplex, the partition function --- and $K$-theoretic
invariants, which discard the Radon--Nikodym derivative, are blind to it. Over a function
field the equilibrium datum happens to be the Weil zeta function itself, which is why the
reconstruction is so complete; but the mechanism is thermodynamic, and asking $KK_\T$ to
supply it repeats an error the programme has already diagnosed three times, in
Remark 4.5 of Chapter~\ref{ch:II}, Proposition 3.1 of Chapter~\ref{ch:III} and Theorem 2.1 of Chapter~\ref{ch:IX}.
\end{remark}

\section{Assessment}

\[
\begin{array}{lll}
\text{Proposal} & \text{Status} & \text{Reference}\\\hline
\partial\cA^{\mathrm{aff}}\ \text{defined for}\ S=\cP_K & \textbf{no: }\cO_{\cP_K}=\F_q & \text{Prop.~1.1}\\
\text{a point at infinity must be fixed} & \text{yes} & \text{Cor.~1.2}\\
\T\text{-actions better behaved than flows} & \textbf{not here} & \text{Prop.~2.1}\\
\text{Gabe--Szab\'o applies after Takai} & \textbf{no: dual algebra stably finite} & \text{Prop.~2.1}\\
\text{Izumi's compact group theory applies} & \textbf{no: action not Rokhlin} & \text{Prop.~2.2}\\
\text{\S7 of Chapter~\ref{ch:XVII} suggestion} & \textbf{withdrawn} & \text{Rem.~2.3}\\
K_0^\T\ \text{encodes}\ P(T) & \textbf{no: only}\ q^{d_S} & \text{Prop.~3.1, Cor.~3.2}\\
K\text{-theory sees}\ |\Jac(\F_q)|=P(1) & \text{yes, one value only} & \text{Rem.~3.3}\\
Z_C\ \text{is the partition function} & \textbf{yes} & \text{Thm.~4.1}\\
\text{complete geometric reconstruction} & \textbf{yes, thermodynamically} & \text{Thm.~4.2}
\end{array}
\]

Two of the three obstructions are specific and one is a retraction of our own suggestion. The
retraction is worth stating plainly: the observation that compact group actions are better
classified than flows is correct about the literature and inapplicable here, because the
standard reduction --- Takai duality to a $\Z$-action --- lands on the stably finite core,
and the compact group theorems that avoid that reduction require a Rokhlin property that
gauge actions do not have. Periodicity of the modular flow is a real structural feature of
the function field case, but it does not buy a classification theorem.

The positive statement is stronger than what was asked for the boundary, and it needs no
classification theory at all. Over a function field the partition function of the Toeplitz
system \emph{is} the Weil zeta function, evaluated at $T=q^{-\beta}$; and since $Z_C$ is
rational, its values on any interval determine everything the Weil conjectures encode. The
system therefore recovers $g$, all $\#C(\F_{q^n})$, the Weil numbers and $|\Jac(C)(\F_q)|$ ---
from thermodynamics.

Three questions. Does the affine boundary with a fixed $\infty$ recover $\deg\infty$, so that
$Z_C$ can be reconstructed from the affine partition function $Z_C(T)(1-T^{\deg\infty})$?
Second, is there any equivariant invariant, short of the full KMS structure, that sees the
$b_n$ --- the graded dimension of the Fock module is the natural candidate, but we do not know
whether it is an invariant of the abstract dynamical system rather than of the chosen
representation. Third, and this is the one we would pursue: two curves with the same zeta
function but non-isomorphic Jacobians exist; by Theorem~\ref{XVIII:thm:reconstruct} their systems
have the same partition function, so does anything in $(\cA,\T,\sigma)$ distinguish them, or
is zeta-equivalence exactly the equivalence realized here?

